% GENERATED FILE. Edit canonical lessons, then run npm run generate:books.
% canonical-source-hash: fnv1a64:34b32a10a946a482
% canonical-lessons: PA-C235-buharna, PA-C235-jharna, PA-C235-pocha, PA-C235-gara, PA-C235-lippna

\chapter{Sweeping the Floor}
\label{ch:pa-a2-sweeping-the-floor}
\hlchaptermodality{235}

\begin{chapteropening}
\textbf{By the end of this chapter:} \emph{I can name the jobs of cleaning a house.}

Complete the last lesson of chapter 235: I can name the jobs of cleaning a house.
\end{chapteropening}

\section[\texorpdfstring{\pa{ਬੁਹਾਰਨਾ} (buhārnā)}{ਬੁਹਾਰਨਾ (buhārnā)}]{\pa{ਬੁਹਾਰਨਾ} (buhārnā) — to sweep}

\label{lesson:PA-C235-buharna}



\begin{quote}
Before the new one: say the Punjabi for to put out, then the Punjabi for a wooden bed.
\end{quote}

\subsection*{You'll want to know: \pa{ਬੁਹਾਰਨਾ}}
\textbf{\pa{ਬੁਹਾਰਨਾ}} — \emph{buhārnā} — \textquotedblleft{}to sweep\textquotedblright{}.

\begin{quote}
\pa{ਮੈਂ} \pa{ਝਾੜੂ} \pa{ਨਾਲ} \pa{ਕਮਰਾ} \pa{ਬੁਹਾਰਦੀ} \pa{ਹਾਂ।}
\end{quote}

\begin{quote}
\emph{maiṁ jhāṛū nāl kamrā buhārdī hāṁ.}
\end{quote}

\begin{quote}
I sweep the room with a broom.
\end{quote}

Past (I …): \pa{ਮੈਂ} \pa{ਕਮਰਾ} \pa{ਬੁਹਾਰਿਆ}, \emph{maiṁ kamrā buhāriā}.

\begin{cousinweb}
\textbf{\pa{ਬੁਹਾਰਨਾ}} (\emph{buhārnā}) is inherited: CDIAL sets it under Sanskrit \emph{bahukāra}, \textquotedblleft{}effecting much\textquotedblright{} (also a late word for a sweeper), and lists Punjabi \emph{buhārnā}, \textquotedblleft{}to sweep\textquotedblright{} (CDIAL 9188).
\end{cousinweb}

\begin{grammarlens}[title={What you've built}]
One more word for this chapter.
\end{grammarlens}

\subsection*{Your turn}
\begin{itemize}
\raggedright
  \item \emph{Say it:} \emph{buhārnā}
  \item \emph{Say it:} \emph{buhārnā}, once more
  \item \emph{Say it:} \emph{maiṁ jhāṛū nāl kamrā buhārdī hāṁ.} — I sweep the room with a broom.
  \item \emph{Recall:} say the Punjabi for to put out, then the Punjabi for a wooden bed, then say \emph{buhārnā} again
\end{itemize}

\begin{tcolorbox}[breakable,colback=teal!4,colframe=teal!35!black,title={Before you move on}]
What does \pa{ਬੁਹਾਰਨਾ} mean? (\textquotedblleft{}To sweep\textquotedblright{}.) Say it once more.
\end{tcolorbox}
\section[\texorpdfstring{\pa{ਝਾੜਨਾ} (jhāṛnā)}{ਝਾੜਨਾ (jhāṛnā)}]{\pa{ਝਾੜਨਾ} (jhāṛnā) — to dust; to shake out}

\label{lesson:PA-C235-jharna}



\begin{quote}
Before the new one: say the Punjabi for a wooden bed, then the Punjabi for to sweep.
\end{quote}

\subsection*{You'll want to know: \pa{ਝਾੜਨਾ}}
\textbf{\pa{ਝਾੜਨਾ}} — \emph{jhāṛnā} — \textquotedblleft{}to dust; to shake out\textquotedblright{}.

\begin{quote}
\pa{ਮੈਂ} \pa{ਕੱਪੜੇ} \pa{ਨਾਲ} \pa{ਅਲਮਾਰੀ} \pa{ਝਾੜਦਾ} \pa{ਹਾਂ।}
\end{quote}

\begin{quote}
\emph{maiṁ kappṛe nāl almārī jhāṛdā hāṁ.}
\end{quote}

\begin{quote}
I dust the cupboard with a cloth.
\end{quote}

Past (I …): \pa{ਮੈਂ} \pa{ਦਰੀ} \pa{ਝਾੜੀ}, \emph{maiṁ darī jhāṛī}.

\begin{cousinweb}
\textbf{\pa{ਝਾੜਨਾ}} (\emph{jhāṛnā}) is inherited from a reconstructed \emph{*jhāṭayati}, \textquotedblleft{}shakes down\textquotedblright{}; CDIAL lists Punjabi \emph{jhāṛṇā}, \textquotedblleft{}to shake off, thresh, sweep\textquotedblright{} (CDIAL 5328).
\end{cousinweb}

\begin{grammarlens}[title={What you've built}]
One more word for this chapter.
\end{grammarlens}

\subsection*{Your turn}
\begin{itemize}
\raggedright
  \item \emph{Say it:} \emph{jhāṛnā}
  \item \emph{Say it:} \emph{jhāṛnā}, once more
  \item \emph{Say it:} \emph{maiṁ kappṛe nāl almārī jhāṛdā hāṁ.} — I dust the cupboard with a cloth.
  \item \emph{Recall:} say the Punjabi for a wooden bed, then the Punjabi for to sweep, then say \emph{jhāṛnā} again
  \item \emph{Write it:} copy \textbf{\pa{ਝਾੜਨਾ}} once, checking the whole word
\end{itemize}

\begin{tcolorbox}[breakable,colback=teal!4,colframe=teal!35!black,title={Before you move on}]
What does \pa{ਝਾੜਨਾ} mean? (\textquotedblleft{}To dust; to shake out\textquotedblright{}.) Say it once more.
\end{tcolorbox}
\section[\texorpdfstring{\pa{ਪੋਚਾ} (pochā)}{ਪੋਚਾ (pochā)}]{\pa{ਪੋਚਾ} (pochā) — a floor cloth (for mopping)}

\label{lesson:PA-C235-pocha}



\begin{quote}
Before the new one: say the Punjabi for to sweep, then the Punjabi for to dust.
\end{quote}

\subsection*{You'll want to know: \pa{ਪੋਚਾ}}
\textbf{\pa{ਪੋਚਾ}} — \emph{pochā} — \textquotedblleft{}a floor cloth (for mopping)\textquotedblright{}.

\begin{quote}
\pa{ਮੈਂ} \pa{ਗਿੱਲੇ} \pa{ਪੋਚੇ} \pa{ਨਾਲ} \pa{ਫ਼ਰਸ਼} \pa{ਸਾਫ਼} \pa{ਕਰਦੀ} \pa{ਹਾਂ।}
\end{quote}

\begin{quote}
\emph{maiṁ gille poche nāl farsh sāf kardī hāṁ.}
\end{quote}

\begin{quote}
I clean the floor with a wet floor cloth.
\end{quote}

A masculine noun.

\begin{cousinweb}
\textbf{\pa{ਪੋਚਾ}} (\emph{pochā}) is inherited from Sanskrit \emph{pōta}, \textquotedblleft{}cloth\textquotedblright{}; CDIAL lists Punjabi \emph{poccā}, \textquotedblleft{}cloth or brush for smearing\textquotedblright{} (CDIAL 8400).
\end{cousinweb}

\begin{grammarlens}[title={What you've built}]
One more word for this chapter.
\end{grammarlens}

\subsection*{Your turn}
\begin{itemize}
\raggedright
  \item \emph{Say it:} \emph{pochā}
  \item \emph{Say it:} \emph{pochā}, once more
  \item \emph{Say it:} \emph{maiṁ gille poche nāl farsh sāf kardī hāṁ.} — I clean the floor with a wet floor cloth.
  \item \emph{Recall:} say the Punjabi for to sweep, then the Punjabi for to dust, then say \emph{pochā} again
  \item \emph{Write it:} copy \textbf{\pa{ਪੋਚਾ}} once, checking the whole word
\end{itemize}

\begin{tcolorbox}[breakable,colback=teal!4,colframe=teal!35!black,title={Before you move on}]
What does \pa{ਪੋਚਾ} mean? (\textquotedblleft{}A floor cloth (for mopping)\textquotedblright{}.) Say it once more.
\end{tcolorbox}
\section[\texorpdfstring{\pa{ਗਾਰਾ} (gārā)}{ਗਾਰਾ (gārā)}]{\pa{ਗਾਰਾ} (gārā) — mud plaster}

\label{lesson:PA-C235-gara}



\begin{quote}
Before the new one: say the Punjabi for to dust, then the Punjabi for a floor cloth.
\end{quote}

\subsection*{You'll want to know: \pa{ਗਾਰਾ}}
\textbf{\pa{ਗਾਰਾ}} — \emph{gārā} — \textquotedblleft{}mud plaster\textquotedblright{}.

\begin{quote}
\pa{ਗਾਰਾ} \pa{ਗਿੱਲਾ} \pa{ਅਤੇ} \pa{ਭਾਰੀ} \pa{ਹੈ।}
\end{quote}

\begin{quote}
\emph{gārā gillā ate bhārī hai.}
\end{quote}

\begin{quote}
The mud plaster is wet and heavy.
\end{quote}

A masculine noun.

\begin{cousinweb}
\textbf{\pa{ਗਾਰਾ}} (\emph{gārā}) is inherited from a reconstructed \emph{*gāra}, \textquotedblleft{}dripping, oozing\textquotedblright{}; CDIAL lists Punjabi \emph{gārā}, \textquotedblleft{}kneaded potter's clay, mud used for mortar\textquotedblright{} (CDIAL 4137).
\end{cousinweb}

\begin{grammarlens}[title={What you've built}]
One more word for this chapter.
\end{grammarlens}

\subsection*{Your turn}
\begin{itemize}
\raggedright
  \item \emph{Say it:} \emph{gārā}
  \item \emph{Say it:} \emph{gārā}, once more
  \item \emph{Say it:} \emph{gārā gillā ate bhārī hai.} — The mud plaster is wet and heavy.
  \item \emph{Recall:} say the Punjabi for to dust, then the Punjabi for a floor cloth, then say \emph{gārā} again
  \item \emph{Write it:} copy \textbf{\pa{ਗਾਰਾ}} once, checking the whole word
\end{itemize}

\begin{tcolorbox}[breakable,colback=teal!4,colframe=teal!35!black,title={Before you move on}]
What does \pa{ਗਾਰਾ} mean? (\textquotedblleft{}Mud plaster\textquotedblright{}.) Say it once more.
\end{tcolorbox}
\section[\texorpdfstring{\pa{ਲਿੱਪਣਾ} (lippṇā)}{ਲਿੱਪਣਾ (lippṇā)}]{\pa{ਲਿੱਪਣਾ} (lippṇā) — to plaster (with mud)}

\label{lesson:PA-C235-lippna}



\begin{quote}
Before the new one: say the Punjabi for a floor cloth, then the Punjabi for mud plaster.
\end{quote}

\subsection*{You'll want to know: \pa{ਲਿੱਪਣਾ}}
\textbf{\pa{ਲਿੱਪਣਾ}} — \emph{lippṇā} — \textquotedblleft{}to plaster (with mud)\textquotedblright{}.

\begin{quote}
\pa{ਦਾਦੀ} \pa{ਗਾਰੇ} \pa{ਨਾਲ} \pa{ਚੁੱਲ੍ਹਾ} \pa{ਲਿੱਪਦੀ} \pa{ਹੈ।}
\end{quote}

\begin{quote}
\emph{dādī gāre nāl chullhā lippdī hai.}
\end{quote}

\begin{quote}
Grandmother plasters the stove with mud.
\end{quote}

Past (I …): \pa{ਮੈਂ} \pa{ਚੁੱਲ੍ਹਾ} \pa{ਲਿੱਪਿਆ}, \emph{maiṁ chullhā lippiā}.

\begin{cousinweb}
\textbf{\pa{ਲਿੱਪਣਾ}} (\emph{lippṇā}) is inherited from Sanskrit \emph{lipyatē}, \textquotedblleft{}is smeared\textquotedblright{}; CDIAL lists Punjabi \emph{lippṇā} among the verbs for plastering a wall (CDIAL 11061).
\end{cousinweb}

\begin{grammarlens}[title={What you've built}]
One more word for this chapter.
\end{grammarlens}

\subsection*{Your turn}
\begin{itemize}
\raggedright
  \item \emph{Say it:} \emph{lippṇā}
  \item \emph{Say it:} \emph{lippṇā}, once more
  \item \emph{Say it:} \emph{dādī gāre nāl chullhā lippdī hai.} — Grandmother plasters the stove with mud.
  \item \emph{Recall:} say the Punjabi for a floor cloth, then the Punjabi for mud plaster, then say \emph{lippṇā} again
  \item \emph{Write it:} copy \textbf{\pa{ਲਿੱਪਣਾ}} once, checking the whole word
\end{itemize}

\begin{tcolorbox}[breakable,colback=teal!4,colframe=teal!35!black,title={Before you move on}]
What does \pa{ਲਿੱਪਣਾ} mean? (\textquotedblleft{}To plaster (with mud)\textquotedblright{}.) Say it once more.
\end{tcolorbox}
